\section{Query complexity, a classical lower bound, and random graphs}\label{sec:complexity}

\subsection{The estimation-free recursion}

\begin{proposition}[Size-biased quantum recursion]
\label{prop:size-biased-algorithm}
Under the hypotheses of \cref{thm:main}, the estimation-free algorithm returns
a verified homogeneous $K$-set with probability at least $1-\eta$ using
\[
 O\!\left(K^2 2^K\log K\log\frac1\eta\right)
\]
edge queries in the worst case.
\end{proposition}

\begin{proof}
A capped sample from a nonempty subset of the $M$-element universe costs
$O(\sqrt M)$ membership tests.  Repeating the block $O(\log(2r))$ times
reduces its failure probability to $O(1/r)$.  Since the membership predicate
at level $i$ costs $O(i+1)$ edge queries, one fixed-cap run uses
\[
 O\!\left(\sqrt M\log(2r)\sum_{i=0}^{r-1}(i+1)\right)
 =O(K^2 2^K\log K)
\]
edge queries.  The survival analysis following
\cref{lem:size-biased-survival} gives a constant success probability for each
run, including capped-search failures and final verification.  Repeating
independent runs $O(\log(1/\eta))$ times proves the stated worst-case bound.
\end{proof}

\subsection{The scale-aware recursion}

\begin{proof}[Proof of \cref{thm:main}]
Set $M=4^{K-1}$ and run the algorithm of \cref{sec:algorithm} on any fixed
set of $M$ input vertices.  The failure allocation in \cref{sec:correctness},
together with \cref{lem:concentration}, implies that all required searches
succeed and all completed estimates are accurate with probability at least
$1-\eta$.  On this event, \cref{lem:survival} guarantees every density
promise used by the sampler and ensures that $S_r$ is nonempty.  The final
output is homogeneous by \cref{lem:output}.  Because the algorithm verifies
the returned set directly, every non-$\bot$ output is valid even outside the
good event.

It remains to bound the number of queries.  By \cref{lem:survival}, a capped
search block at level $i$ uses $O(2^{i/2})$ membership tests.  Computing and
uncomputing membership and querying the sampled edge costs $O(i+1)$ edge
queries.  Define
\[
 b_i=(i+1)2^{i/2}.
\]
The batch cap in \cref{lem:capped-search} uses
$O(m_i+\log(r/\eta))=O(m_i)$ blocks.  By \cref{eq:sample-count}, the total
cost of the sampling batches is therefore
\begin{equation}\label{eq:query-sum}
 O\!\left(
 \log\frac r\eta
 \sum_{i=0}^{r-1}b_i\varepsilon_i^{-2}
 \right).
\end{equation}

At most six indices share any value of $d_i$, and hence
\[
 Z=\sum_i2^{-d_i}<6\sum_{d\ge0}2^{-d}=12.
\]
Using \cref{eq:error-schedule},
\[
 \sum_{i=0}^{r-1}b_i\varepsilon_i^{-2}
 =O\!\left(
 \sum_{i=0}^{r-1}(i+1)2^{i/2}4^{d_i}
 \right).
\]
Write $h=r-1-i$.  Since
$4^{\lceil h/6\rceil}\le4\cdot2^{h/3}$, the last display is at most
\[
 O\!\left(
 2^{r/2}\sum_{h=0}^{r-1}(r-h)2^{-h/6}
 \right)
 =O(r2^{r/2}).
\]
The amplified pivot searches and the final search cost
\[
 O\!\left(\log\frac r\eta\sum_{i=0}^{r}b_i\right)
 =O\!\left(r2^{r/2}\log\frac r\eta\right),
\]
and the $O(K^2)$ output verification is smaller.  Substituting $r=2K-3$
into \cref{eq:query-sum} gives
\[
 O\!\left(2^K K\log\frac K\eta\right),
\]
as claimed.
\end{proof}

Each search begins in the uniform state on the fixed power-of-two universe
$[M]$.  The recorded pivot labels and colours define the reversible membership
predicate for $S_i$, so the algorithm requires no materialized representation
of any candidate set.

\subsection{A randomized classical lower bound}
\label{sec:classical-lower}

We complement the upper bound with a randomized classical lower bound.  It
is the random-Painter estimate of Conlon, Fox, Grinshpun, and He
\cite[Theorem~1]{ConlonFoxGrinshpunHe2019} for online Ramsey numbers,
transported to the query model.  Because that estimate is a statement about
uniformly random colourings, the bound holds on the uniform distribution
$G(N,1/2)$ and not only in the worst case.  \Cref{app:classical-lower}
gives a self-contained proof.

\begin{proposition}[Randomized classical lower bound]
\label{prop:classical-lower}
Let $K\ge8$ and let $T=\lfloor2^{(2-\sqrt2)K-2}\rfloor$.  For every
$N\ge K$, every randomized algorithm that makes at most $T$ edge queries
outputs a homogeneous $K$-set of $G\sim G(N,1/2)$ with probability at most
$5/8$.
Consequently, for $M=4^{K-1}$, every randomized edge-query algorithm that
finds a homogeneous $K$-set with worst-case probability at least $2/3$ uses
\[
 \Omega\!\left(2^{(2-\sqrt2)K}\right)
 =\Omega\!\left(M^{1-1/\sqrt2}\right)
\]
queries.
\end{proposition}

The exponent $1-1/\sqrt2\approx0.2929$ improves the $M^{1/4}$ black-box
lower bound of \cite{ImpagliazzoNaor1988,KomargodskiNaorYogev2019}.  Jain
et al.\ \cite[Section~III]{JainLiRobereXun2024} already point to
\cite{ConlonFoxGrinshpunHe2019} for the best known deterministic lower
bound; \cref{prop:classical-lower} records the randomized and
distributional form that the same analysis yields.  It leaves the
randomized worst-case complexity between $\Omega(M^{1-1/\sqrt2})$ and the
$O(M)$ of the explicit recursion, and it does not separate
\cref{thm:main} from every classical algorithm in the worst case.

\subsection{A quantum speedup on random graphs}
\label{sec:random-graphs}

On the distribution $G(M,1/2)$ of \cref{prop:classical-lower}, a greedy
quantum neighbourhood search is much cheaper than the worst-case algorithm
of \cref{thm:main}: the common neighbourhood of every small clique in a
random graph has essentially its expected size, so no majority has to be
estimated and the recursion has depth $K$ rather than $2K-3$.

\begin{proposition}[Quantum greedy search on random graphs]
\label{prop:random-graph}
Let $K\ge2$, $M=4^{K-1}$, and $0<\eta<1/2$, and let $G\sim G(M,1/2)$.
There is a quantum algorithm which makes
\[
 O\!\left(K2^{K/2}\log\frac K\eta\right)
\]
edge queries on every input, never outputs a set that is not a $K$-clique,
and outputs a $K$-clique with probability at least $1-\eta-\pi_K$ over $G$
and its internal randomness, where $\pi_K=2M^{K-1}e^{-2^{K-2}/12}$
satisfies $\pi_K<e^{-100}$ for all $K\ge14$.
\end{proposition}

The proof, given in \cref{app:random-graph}, shows that with probability
at least $1-\pi_K$ every $j$-set with $j<K$ has at least $M2^{-j-1}$ common
neighbours, and then runs the capped search of \cref{lem:capped-search}
with density parameter $2^{-j-1}$ at level $j$.  Combined with the
distributional form of \cref{prop:classical-lower}, this gives a provable
quantum speedup for Ramsey search on random graphs.

\begin{corollary}[Separation on random graphs at the Ramsey scale]
\label{cor:separation}
Let $n\ge14$, $N=2^n$, $K=\lfloor n/2\rfloor+1$, $0<\eta<1/2$, and
$G\sim G(N,1/2)$.
\begin{enumerate}
\item A quantum algorithm making
$O(N^{1/4}\log N\log(\log N/\eta))$ edge queries outputs a homogeneous
$K$-set of $G$ with probability at least $1-\eta-\pi_K$.
\item Every randomized algorithm making at most
$\lfloor2^{(2-\sqrt2)K-2}\rfloor=\Omega(N^{1-1/\sqrt2})$ edge queries
outputs a homogeneous $K$-set of $G$ with probability at most $5/8$.
\end{enumerate}
Since $1/4<1-1/\sqrt2$, the bounded-error quantum query complexity of the
Ramsey relation on $G(N,1/2)$ is polynomially smaller than its bounded-error
randomized query complexity.
\end{corollary}

\begin{proof}
The subgraph induced on the first $M=4^{K-1}\le N$ vertices is distributed
as $G(M,1/2)$, so the first part follows from \cref{prop:random-graph},
using $2^{K/2}\le2^{n/4+1/2}$ and $K=O(\log N)$.  The second part is the
distributional statement of \cref{prop:classical-lower}, which applies
because $K\ge8$.  For the final sentence, take $\eta=1/4$ and $K\ge14$, so
that the quantum algorithm succeeds with probability greater than $2/3$
while $5/8<2/3$.
\end{proof}

The size-biased recursion also extends to every fixed number of edge colours.
The statement and proof appear in \cref{app:multicolour}.
